% =============================================================================
%  Question 1 --- Lecture 2: Reification in property graphs
%  Source: docs/course_extracts/lecture2/06-exam-and-readings.md
% =============================================================================

\begin{question}{1}{2}{Reification in property graphs}

\questionbrief
  {%
    In this lecture we have discussed the notion of reification for RDF.
    Formulate a notion of reification in the world of property graphs and
    illustrate how reification in property graphs would work in practice using
    two small examples. One of those examples should be
    \emph{``The detective supposes that the butler killed the gardener''}.
    Since we are dealing with property graphs, make sure to include
    vertex/edge labels and key-value pairs in your proposal.%
  }
  {%
    No specific paper is required by the prompt. Useful background:
    Hogan et al., \emph{Knowledge graphs} (chs.~1--2);
    Angles \& Guti\'errez, \emph{Survey of graph database models} (ACM CSUR, 2008);
    Guarino, Sales \& Guizzardi, \emph{Reification and Truthmaking Patterns}
    (ER 2018).%
  }
  {%
    70\% completeness, correctness, and elegance of the proposal;
    30\% writing quality/clarity/completeness.
    Minimum length: \textbf{2 A4 full pages}.%
  }

% --- Your answer below -------------------------------------------------------

% TODO: define your notion of property-graph reification (motivation +
% formal/semi-formal definition), then provide the two required examples
% with explicit vertex/edge labels and key-value pairs.

\end{question}
